% Alignment foundations, scalable search, and the limits that motivate BioSeqDB.

\section{Alignment objectives and boundary conditions}
\label{sec:alignment-model}
\label{subsec:alignment-modes}

Sequence search and sequence alignment answer different questions. Search
identifies locations that may contain a useful correspondence. Alignment then
describes one correspondence at symbol resolution. This distinction matters at
genome scale. Exact dynamic programming can optimize a stated pairwise
problem, but applying it to every reference interval would be too costly.
Practical search methods first reduce the candidate set, then align only the
retained regions.

Let \(X=x_1\ldots x_n\) and \(Y=y_1\ldots y_m\) be the sequence portions
consumed by one traceback. Under a global objective they are the complete
inputs. Other boundary modes may select substrings or leave terminal flanks
unaligned. An alignment is a sequence of columns
\[
  A=((a_1,b_1),\ldots,(a_L,b_L)),
  \qquad
  (a_r,b_r)\in(\Sigma\cup\{-\})^2\setminus\{(-,-)\},
\]
whose first and second rows recover \(X\) and \(Y\) after gaps are removed. A
column with two symbols is a match or substitution. A symbol opposite a gap is
an insertion or deletion relative to the other sequence. The alignment is an
optimized correspondence under a scoring model. Its score alone does not
establish evolutionary homology.

For substitution score \(s(a,b)\) and gap cost \(g(\ell)\), the score is
\[
  S(A)=\sum_{\text{symbol pairs }r}s(a_r,b_r)
       -\sum_{\text{gap runs }u}g(|u|).
\]
This chapter uses the affine convention
\begin{equation}
  g(\ell)=q+\ell e,
  \qquad \ell\geq 1,
  \label{eq:affine-gap-convention}
\end{equation}
where \(q\) is the opening cost and \(e\) is the extension cost. A gap of length
one therefore costs \(q+e\). Sources that write \(q+(\ell-1)e\) use a different
parameter convention. Affine scoring charges one opening cost per gap run and
can be evaluated with a fixed number of states per matrix cell
\cite{gotoh_1982_affine}.

The scoring recurrence does not fully define the problem. Initialization,
restart rules, and permitted traceback endpoints determine which sequence
ends must be accounted for. Global alignment consumes both complete sequences
\cite{needleman_wunsch_1970}. Local alignment may restart at zero and selects
the best pair of substrings \cite{smith_waterman_1981}. Semi-global and overlap
objectives consume selected ends while allowing specified flanks to remain
unaligned. Extension begins from a fixed anchor and chooses an endpoint under
a stopping rule. Clipping records a terminal segment as unaligned rather than
charging it as an ordinary gap. Local-extension fixes the start but permits an
interior endpoint. Figure~\ref{fig:alignment-boundary-modes} shows schematic
examples of these modes.

\definecolor{alignmentamber}{HTML}{EFB84C}
\definecolor{alignmentcream}{HTML}{FBF6D5}
\definecolor{alignmentgray}{HTML}{F3F3F3}
\definecolor{alignmentlight}{HTML}{F9F9F9}

\begin{figure}[tbp]
  \centering
  \begin{tikzpicture}[
      x=1cm,
      y=1cm,
      >=Latex,
      font=\footnotesize,
      panel title/.style={font=\small,align=center},
      matrix border/.style={draw=black,line width=0.45pt},
      penalty edge/.style={draw=black,line width=1.35pt,line cap=round},
      path line/.style={draw=alignmentamber,line width=1.15pt,
                        line cap=round,line join=round},
      query bar/.style={draw=black,dashed,fill=alignmentcream,
                        line width=0.55pt}
    ]

    % Global
    \begin{scope}[shift={(0,3.55)}]
      \node[panel title] at (1.8,2.78) {Global};
      \draw[matrix border] (1.05,1.02) rectangle (2.55,2.52);
      \draw[path line] (1.05,2.52) -- (1.40,2.17) -- (1.78,2.17)
        -- (1.78,1.78) -- (2.14,1.42) -- (2.55,1.02);
      \fill[black] (1.05,2.52) circle[radius=1.7pt];
      \fill[black] (2.55,1.02) circle[radius=1.7pt];
      \draw[-{Latex[length=1.8mm]},line width=0.55pt]
        (0.32,0.28) -- (3.28,0.28);
      \fill[alignmentcream] (0.38,0.52) rectangle (3.18,0.78);
      \draw[line width=0.55pt] (0.38,0.78) -- (3.18,0.78);
      \draw[dashed,line width=0.55pt] (0.38,0.52) -- (3.18,0.52);
    \end{scope}

    % Ends-free
    \begin{scope}[shift={(3.8,3.55)}]
      \node[panel title] at (1.8,2.78) {Ends-free};
      \draw[matrix border] (0.30,1.15) rectangle (3.30,2.40);
      \draw[penalty edge] (0.30,1.90) -- (0.30,2.40) -- (1.50,2.40);
      \draw[penalty edge] (2.40,1.15) -- (3.30,1.15) -- (3.30,1.53);
      \draw[path line] (1.48,2.40) -- (1.75,2.14) -- (2.05,2.14)
        -- (2.05,1.86) -- (2.48,1.43) -- (2.86,1.15);
      \draw[-{Latex[length=1.8mm]},line width=0.55pt]
        (0.30,0.28) -- (3.30,0.28);
      \draw[query bar] (0.82,0.52) rectangle (2.78,0.78);
    \end{scope}

    % Semi-global
    \begin{scope}[shift={(7.6,3.55)}]
      \node[panel title] at (1.8,2.78) {Semi-global};
      \draw[matrix border] (0.30,1.15) rectangle (3.30,2.40);
      \draw[penalty edge] (0.30,2.40) -- (3.30,2.40);
      \draw[penalty edge] (0.30,1.15) -- (3.30,1.15);
      \draw[path line] (1.20,2.40) -- (1.52,2.10) -- (1.82,2.10)
        -- (1.82,1.82) -- (2.34,1.32) -- (3.02,1.15);
      \draw[-{Latex[length=1.8mm]},line width=0.55pt]
        (0.30,0.28) -- (3.30,0.28);
      \fill[alignmentcream] (1.18,0.52) rectangle (3.02,0.78);
      \draw[dashed,-{Latex[length=1.5mm]},line width=0.55pt]
        (1.18,0.52) -- (2.85,0.52)
        -- (3.08,0.34);
      \draw[-{Latex[length=1.7mm]},line width=0.55pt]
        (1.18,0.78) -- (3.08,0.78);
    \end{scope}

    % Local
    \begin{scope}[shift={(11.4,3.55)}]
      \node[panel title] at (1.8,2.78) {Local};
      \draw[matrix border,fill=alignmentgray]
        (1.05,1.02) rectangle (2.55,2.52);
      \draw[path line] (1.40,2.28) -- (1.62,2.06) -- (1.90,2.06)
        -- (1.90,1.78) -- (2.18,1.50);
      \fill[alignmentamber] (1.40,2.28) circle[radius=1.8pt];
      \fill[alignmentamber] (2.18,1.50) circle[radius=1.8pt];
      \draw[-{Latex[length=1.8mm]},line width=0.55pt]
        (0.30,0.28) -- (3.30,0.28);
      \draw[query bar] (1.03,0.52) rectangle (2.57,0.78);
      \draw[dashed,-{Latex[length=1.5mm]}]
        (1.03,0.64) -- (0.58,0.35);
      \draw[dashed,-{Latex[length=1.5mm]}]
        (2.57,0.64) -- (3.02,0.35);
    \end{scope}

    % Extension
    \begin{scope}
      \node[panel title] at (1.8,2.78) {Extension};
      \draw[matrix border] (0.30,1.15) rectangle (3.30,2.40);
      \draw[penalty edge] (0.30,1.15) -- (3.30,1.15) -- (3.30,2.40);
      \draw[path line] (0.30,2.40) -- (0.66,2.08) -- (1.02,2.08)
        -- (1.02,1.76) -- (2.10,1.15);
      \fill[black] (0.30,2.40) circle[radius=1.7pt];
      \draw[-{Latex[length=1.8mm]},line width=0.55pt]
        (0.30,0.28) -- (3.30,0.28);
      \fill[alignmentcream] (0.30,0.52) rectangle (2.05,0.78);
      \draw[-{Latex[length=1.7mm]},line width=0.55pt]
        (0.30,0.65) -- (2.15,0.65);
      \draw[dashed] (2.78,0.55) -- (3.04,0.34);
    \end{scope}

    % Overlap
    \begin{scope}[shift={(3.8,0)}]
      \node[panel title] at (1.8,2.78) {Overlap};
      \draw[matrix border] (0.30,1.15) rectangle (3.30,2.40);
      \draw[penalty edge] (0.30,2.40) -- (0.30,1.15) -- (3.30,1.15);
      \draw[path line] (0.30,2.02) -- (0.68,1.68) -- (1.05,1.68)
        -- (1.43,1.34) -- (2.25,1.15);
      \draw[-{Latex[length=1.8mm]},line width=0.55pt]
        (0.30,0.28) -- (3.30,0.28);
      \draw[query bar] (0.54,0.52) rectangle (2.38,0.78);
      \draw[dashed,-{Latex[length=1.5mm]}]
        (0.54,0.64) -- (0.17,0.35);
    \end{scope}

    % Clipping
    \begin{scope}[shift={(7.6,0)}]
      \node[panel title] at (1.8,2.78) {Clipping};
      \draw[fill=alignmentlight,draw=black,line width=1.35pt]
        (0.30,1.15) rectangle (3.30,2.40);
      \draw[path line] (0.30,2.40) -- (0.70,2.08) -- (1.14,2.08)
        -- (1.14,1.75) -- (2.68,1.58);
      \draw[path line,dashed] (2.68,1.58) -- (2.68,1.15);
      \fill[black] (0.30,2.40) circle[radius=1.7pt];
      \fill[black] (3.30,1.15) circle[radius=1.7pt];
      \fill[alignmentamber] (2.68,1.58) circle[radius=1.8pt];
      \draw[-{Latex[length=1.8mm]},line width=0.55pt]
        (0.30,0.28) -- (3.30,0.28);
      \draw[query bar] (0.80,0.52) rectangle (2.60,0.78);
      \draw[-{Latex[length=1.6mm]},line width=0.55pt]
        (2.60,0.65) -- (2.72,0.65) -- (2.92,0.38);
    \end{scope}

    % Local extension
    \begin{scope}[shift={(11.4,0)}]
      \node[panel title] at (1.8,2.78) {Local-extension};
      \draw[matrix border,fill=alignmentlight]
        (1.05,1.02) rectangle (2.55,2.52);
      \draw[path line] (1.05,2.52) -- (1.36,2.20) -- (1.68,2.20)
        -- (1.68,1.88) -- (2.05,1.50);
      \fill[black] (1.05,2.52) circle[radius=1.7pt];
      \fill[alignmentamber] (2.05,1.50) circle[radius=1.8pt];
      \draw[-{Latex[length=1.8mm]},line width=0.55pt]
        (0.30,0.28) -- (3.30,0.28);
      \draw[query bar] (0.30,0.52) rectangle (2.05,0.78);
      \draw[dashed,-{Latex[length=1.5mm]}]
        (2.05,0.64) -- (2.62,0.35);
    \end{scope}

    \draw[path line] (0.35,-0.20) -- (0.75,-0.20);
    \node[anchor=west] at (0.84,-0.20) {alignment path};
    \fill[black] (3.70,-0.20) circle[radius=1.7pt];
    \node[anchor=west] at (3.82,-0.20) {fixed endpoint};
    \fill[alignmentamber] (6.75,-0.20) circle[radius=1.8pt];
    \node[anchor=west] at (6.88,-0.20) {selected free endpoint};
    \draw[penalty edge] (10.60,-0.20) -- (11.02,-0.20);
    \node[anchor=west] at (11.12,-0.20) {penalized border};
  \end{tikzpicture}
  \caption{Eight alignment boundary modes. Amber paths show possible
    tracebacks rather than scores computed for one sequence pair. Black dots
    mark fixed endpoints, amber dots mark selected free endpoints, and bold
    border segments mark penalized or anchored boundaries. In each lower
    schematic, the solid arrow is the reference and the cream dashed track is
    the query. These panels summarize boundary patterns and do not replace a
    library's precise initialization and endpoint rules.}
  \label{fig:alignment-boundary-modes}
\end{figure}

The mode name is not sufficient when results are compared across libraries.
The axis assignment and boundary initialization must be stated. The permitted
traceback start and endpoint must also be given. Local alignment further needs
the zero-restart rule, while extension may use an early stopping criterion.

\section{Exact dynamic programming}
\label{sec:exact-dp}

Let \(H_{i,j}\) be the best score after consuming \(i\) symbols of \(X\) and
\(j\) symbols of \(Y\). With a linear gap cost \(d\), a global recurrence is
\begin{equation}
  H_{i,j}=\max\bigl\{
    H_{i-1,j-1}+s(x_i,y_j),
    H_{i-1,j}-d,
    H_{i,j-1}-d
  \bigr\}.
  \label{eq:linear-gap-dp}
\end{equation}
Its boundary row and column accumulate gap costs. A local recurrence adds zero
as a candidate. Free-end objectives change the boundary values and the cells
from which traceback may begin.

Affine scoring requires a main state and one state for each gap orientation.
Let \(I_{i,j}\) and \(D_{i,j}\) end in gaps in \(X\) and \(Y\), respectively,
and let \(M_{i,j}\) end in a symbol pair. Under
Equation~\ref{eq:affine-gap-convention}, one formulation is
\begin{equation}
\begin{aligned}
  M_{i,j} &= H_{i-1,j-1}+s(x_i,y_j),\\
  I_{i,j} &= \max\{H_{i,j-1}-q-e,\ I_{i,j-1}-e\},\\
  D_{i,j} &= \max\{H_{i-1,j}-q-e,\ D_{i-1,j}-e\},\\
  H_{i,j} &= \max\{M_{i,j},I_{i,j},D_{i,j}\}.
\end{aligned}
\label{eq:affine-dp}
\end{equation}
Using \(H\) to open a gap permits a transition from any previous state. Local
alignment also includes zero in the final maximum. Sentinel values exclude
states that cannot occur at a selected boundary.

\begin{algorithm}[tb]
  \caption{Affine-gap alignment with an explicit boundary contract}
  \label{alg:affine-dp}
  \begin{algorithmic}[1]
    \Require sequences \(X[1\ldots n]\), \(Y[1\ldots m]\), scores \(s,q,e\),
      and a boundary contract \(B\)
    \Ensure score, endpoints, and an optional CIGAR string
    \State initialize \(H\), \(I\), and \(D\) according to \(B\)
    \State \textbf{for} \(i\gets 1\) to \(n\) \textbf{do}
      \State \hspace{\algorithmicindent}\textbf{for} \(j\gets 1\) to \(m\)
        \textbf{do}
        \State update \(M_{i,j}\), \(I_{i,j}\), \(D_{i,j}\), and \(H_{i,j}\)
          with Equation~\ref{eq:affine-dp}
        \State store the predecessor of each winning state when traceback is
          requested
      \State \hspace{\algorithmicindent}\textbf{end for}
    \State \textbf{end for}
    \State select the endpoint permitted by \(B\)
    \State follow predecessors to an allowed start and run-length encode the
      operations
  \end{algorithmic}
\end{algorithm}

Traceback yields an edit path that can be encoded as a CIGAR string. This is
the standard compact representation of gapped alignments
\cite{sam_format_specification}.
Unbanded matrix filling costs \(O(nm)\) time. A score needs only a small number
of rows, while direct traceback normally retains predecessor information for
the evaluated cells.

A two-piece affine model uses
\[
  g_2(\ell)=\min(q_1+\ell e_1,\ q_2+\ell e_2)
\]
and adds a second state for each gap orientation. This permits different cost
slopes for short and long gaps. The minimap2-derived route can use this model,
while an adaptive block route uses one affine component.

Large alignments require a smaller evaluated region. A band retains cells near
an expected diagonal. Its cost is proportional to the band width and sequence
length, but it changes the result if the best path leaves the band. X-drop
stops an extension when its score falls sufficiently below the best score seen
so far. Z-drop also considers displacement from the best diagonal, which makes
it less likely that one long gap is treated as an ordinary score collapse.
Adaptive block alignment follows a promising rectangular region rather than a
fixed band. It can also miss the unconstrained optimum
\cite{block_aligner_paper}.

Vector execution changes how cells are evaluated, not the intended objective.
Equivalent scalar and vector routes still need the same precision, boundary
values, cell set, stopping rule, and tie policy. The Suzuki and Kasahara
difference formulation stores bounded differences between neighboring scores,
which supports narrow integer lanes and forms the basis of KSW2
\cite{suzuki_kasahara_2018}. BioSeqDB uses portable SIMD for selected
alignment routes and configurations.

\section{Search by seeding, chaining, and extension}
\label{sec:seed-chain-extend}

Applying local dynamic programming between a query of length \(m\) and a
reference of length \(N\) costs \(O(Nm)\). Most reference positions cannot
belong to a useful result. Long-read mappers therefore organize search into
three stages \cite{minimap2_paper}.

\begin{enumerate}
  \item \textbf{Seeding} selects short query substrings and finds compatible
    reference occurrences.
  \item \textbf{Chaining} selects ordered, approximately collinear hits and
    turns them into candidate regions.
  \item \textbf{Extension} performs base-level dynamic programming between and
    around retained anchors, then returns coordinates, a score, and optional
    traceback data.
\end{enumerate}

An exact \(k\)-mer seed is a length-\(k\) word shared by the query and target.
Indexing every \(k\)-mer produces many redundant neighboring hits. A minimizer
sketch selects a representative \(k\)-mer from each window of \(w\) consecutive
\(k\)-mers, commonly the minimum under a hash order
\cite{roberts_minimizers_2004}. Overlapping windows often retain the same
representative, which is emitted once. Canonical nucleotide seeds use one key
for a \(k\)-mer and its reverse complement while retaining orientation.
Occurrence filtering removes or down-weights very frequent keys because they
carry little location information.

A lookup produces anchors. Write anchor \(a_i\) as a reference coordinate
\(r_i\), query coordinate \(q_i\), orientation \(o_i\), and span \(\ell_i\).
For reverse-strand hits, transform \(q_i\) into the coordinate order of the
oriented query. Denote the resulting coordinate by \(\tilde q_i\). Anchors from
one ungapped correspondence then lie near a common diagonal
\(r_i-\tilde q_i\). Insertions and deletions shift that diagonal. Chaining is
therefore a sparse path problem in the oriented reference-query plane.

After sorting anchors by the canonical coordinate order, a representative
score is
\begin{equation}
  f(i)=\max\left\{
    \ell_i,
    \max_{j<i}\bigl(f(j)+\rho(a_j,a_i)-\gamma(a_j,a_i)\bigr)
  \right\},
  \label{eq:chaining-score}
\end{equation}
where \(\rho\) rewards newly covered matching span and \(\gamma\) penalizes
large gaps and diagonal changes. A predecessor is admitted only when its
orientation, coordinate order, distance, and bandwidth satisfy the chaining
policy. Backtracking predecessor indices from high-scoring endpoints recovers
chains.

Testing every predecessor costs \(O(A^2)\) for \(A\) anchors. Minimap2 limits
the scan to recent admissible anchors and may stop after a configured number of
unproductive candidates. These rules reduce work in practice, but the optimum
is then defined over the examined predecessor set. Range-minimum-query methods
accelerate a related long-gap formulation. SIMD and GPU methods instead
evaluate groups of predecessor candidates together. These approaches may
share a scoring function without sharing the same candidate policy.

\begin{algorithm}[tb]
  \caption{Conceptual seed, chain, and extend search}
  \label{alg:seed-chain-extend}
  \begin{algorithmic}[1]
    \Require query \(Q\), compatible target index \(I\), and search policy \(P\)
    \Ensure ranked alignments
    \State sketch \(Q\) and retrieve compatible seed occurrences from \(I\)
    \State form oriented anchors and suppress uninformative frequent keys
    \State sort anchors in canonical coordinate order
    \State \textbf{for each} anchor endpoint \(i\) \textbf{do}
      \State score the predecessor candidates admitted by \(P\)
      \State store the best score and predecessor for \(i\)
    \State \textbf{end for}
    \State backtrack high-scoring chains and derive candidate regions
    \State \textbf{for each} retained region \textbf{do}
      \State align inter-anchor gaps and extend the two outer regions
      \State compose endpoints, score, and optional CIGAR data
    \State \textbf{end for}
    \State rank competing mappings and compute mapping quality
  \end{algorithmic}
\end{algorithm}

The final alignment is assembled from related boundary problems. Inter-anchor
segments usually require both anchor ends. Outer regions are extension
problems and may terminate under X-drop or Z-drop. Mapping quality (MAPQ) then
summarizes one mapping relative to alternatives and repetitive evidence. It is
not the score of one dynamic-programming matrix.

BioSeqDB retains the minimap2 stages but accepts alphabet-aware compact ranges,
manages persistent target indexes, batches database rows, and returns typed
alignment fields. Individual stages can select scalar, portable SIMD, or GPU
execution when their configuration supports it. Backend selection does not
change the meaning of an alignment column.

\section{Profile database search}
\label{sec:profile-database-search}

Protein profile search uses a different scoring model but applies the same
principle of progressive candidate reduction. In MMseqs2, two successive
indexed hits for one target with equal diagonals activate ungapped scoring.
Only a smaller candidate set reaches profile alignment
\cite{mmseqs2_2017}. Figure~\ref{fig:profile-double-hit} shows this funnel and
the loop order behind the diagonal check.

\definecolor{profilequery}{HTML}{EE3B3B}
\definecolor{profiletarget}{HTML}{2A2AE6}
\definecolor{profileposition}{HTML}{F08A22}
\definecolor{profilehit}{HTML}{C21FA0}
\definecolor{profilematch}{HTML}{1EA01E}

\begin{figure}[tbp]
  \centering
  \begin{tikzpicture}[
      x=1cm,
      y=1cm,
      >=Latex,
      font=\sffamily\footnotesize,
      stage caption/.style={align=center,text width=3.2cm},
      plot/.style={draw=black,fill=white,line width=0.55pt},
      neutral arrow/.style={draw=gray!55,line width=1.7mm,
                            -{Latex[length=3.5mm,width=4mm]}},
      loop badge/.style={circle,minimum size=5.5mm,inner sep=0pt,
                         text=white,font=\small\bfseries}
    ]

    \node[font=\small\bfseries,anchor=west] at (0,13.55)
      {(a) Candidate funnel};

    % Fast k-mer match deck
    \foreach \o in {0.30,0.15} {
      \begin{scope}[shift={(\o,\o)}]
        \draw[draw=gray!55,fill=white] (0.48,10.25) rectangle (3.08,12.72);
      \end{scope}
    }
    \draw[plot] (0.48,10.25) rectangle (3.08,12.72);
    \draw[gray!55] (0.82,10.55) -- (2.78,12.38);
    \draw[gray!55] (0.82,10.72) -- (2.67,12.45);
    \draw[fill=gray!30,rounded corners=0.6mm]
      (1.02,10.86) rectangle (1.29,10.99);
    \draw[fill=black,rounded corners=0.6mm]
      (1.48,11.25) rectangle (1.79,11.39);
    \draw[fill=black,rounded corners=0.6mm]
      (2.04,11.77) rectangle (2.35,11.91);
    \draw[fill=gray!30,rounded corners=0.6mm]
      (2.52,11.52) rectangle (2.79,11.65);
    \node[text=profiletarget,rotate=90,font=\bfseries] at (0.18,11.48)
      {target profile};
    \node[text=profilequery,font=\bfseries] at (1.78,10.02) {query};
    \node[align=center,text width=2.9cm] at (1.78,13.02)
      {many target profiles};
    \node[stage caption] at (1.78,9.55) {Fast \(k\)-mer match};
    \node[align=center] at (1.82,12.26) {two hits on\\one diagonal};
    \draw[->] (1.62,11.98) -- (1.63,11.42);

    \draw[neutral arrow] (3.55,11.48) -- (5.55,11.48);
    \node[align=center,text width=2.1cm] at (4.55,12.16)
      {double hit on\\the same diagonal};

    % Ungapped alignment deck
    \foreach \o in {0.24,0.12} {
      \begin{scope}[shift={(\o,\o)}]
        \draw[draw=gray!55,fill=white] (6.00,10.25) rectangle (8.60,12.72);
      \end{scope}
    }
    \draw[plot] (6.00,10.25) rectangle (8.60,12.72);
    \draw[black,line width=1.25pt] (6.38,10.65) -- (7.24,11.45);
    \draw[black,line width=1.25pt] (7.38,11.58) -- (8.25,12.38);
    \node[align=center,text width=2.9cm] at (7.30,13.02)
      {fewer target profiles};
    \node[text=profilequery,font=\bfseries] at (7.30,10.02) {query};
    \node[stage caption] at (7.30,9.55)
      {Ungapped alignment of target consensus};

    \draw[neutral arrow] (9.05,11.48) -- (11.05,11.48);
    \node[align=center,text width=2.0cm] at (10.05,12.16)
      {transpose\\shortlist};

    % Final profile alignment
    \draw[plot] (11.48,10.25) rectangle (14.08,12.72);
    \draw[black,line width=1.25pt] (11.85,10.65) -- (12.65,11.38)
      -- (12.92,11.38) -- (13.72,12.15);
    \node[align=center,text width=2.9cm] at (12.78,13.02)
      {final candidates};
    \node[text=profilequery,font=\bfseries] at (12.78,10.02) {query};
    \node[stage caption] at (12.78,9.55)
      {Target-profile alignment};

    \draw[gray!55,dashed] (0,8.92) -- (15,8.92);
    \node[font=\small\bfseries,anchor=west] at (0,8.55)
      {(b) Nested-loop double-hit search};

    \node[draw=profilequery,cylinder,shape border rotate=90,aspect=0.28,
          minimum width=2.55cm,minimum height=1.05cm,align=center,
          text=profilequery] (queries) at (1.48,6.82)
      {query 1\\\(\vdots\)\\query \(x\)};
    \node[loop badge,fill=profilequery] at (0.24,7.72) {1};
    \node[anchor=west,text=profilequery,font=\bfseries] at (0.58,7.72)
      {queries};
    \node[draw=profileposition,dashed,rounded corners=2pt,
          minimum width=2.9cm,minimum height=0.78cm,align=center,
          text=profilequery] (qmer) at (1.48,4.92)
      {\textcolor{profileposition}{position \(i=151\)}\\
       \textcolor{profilequery}{\texttt{MHVWAVR}}};
    \draw[->,draw=profilequery,line width=0.8pt] (queries) -- (qmer);
    \node[loop badge,fill=profileposition] at (0.24,5.68) {2};
    \node[anchor=west,text=profileposition,font=\bfseries] at (0.58,5.68)
      {positions};

    \draw[draw=profiletarget,rounded corners=2pt,line width=0.8pt]
      (3.55,3.78) rectangle (9.15,7.06);
    \node[text=profiletarget,font=\small\bfseries] at (6.35,6.55)
      {Precomputed profile index};
    \node[draw=profiletarget,cylinder,shape border rotate=90,aspect=0.28,
          minimum width=2.55cm,minimum height=0.82cm,align=center,
          text=profiletarget] (targets) at (6.35,7.65)
      {target profiles};
    \draw[->,draw=profiletarget,line width=0.8pt]
      (targets.south) -- (6.35,7.06);
    \draw[gray!50] (6.30,4.03) -- (6.30,6.20);
    \node[align=center,text width=2.35cm] at (4.88,5.70)
      {similar \(k\)-mers\\above threshold};
    \node[draw=profileposition,rounded corners=1pt] at (4.88,4.65)
      {\texttt{MHWVREA}};
    \node[text=profilehit,font=\bfseries] at (7.46,5.70) {index hits};
    \node[loop badge,fill=profilehit] at (8.67,5.70) {3};
    \node[align=left] at (7.72,4.60)
      {\((23,431)\)\\\textcolor{profilehit}{\((972,149)\)}\\
       \((1522,217)\)};
    \draw[->,draw=profileposition,line width=0.8pt]
      (qmer.east) -- node[above] {lookup} (4.35,4.65);
    \draw[plot] (10.15,3.78) rectangle (14.55,7.08);
    \node[text=profilehit,rotate=90,font=\bfseries] at (9.87,5.43)
      {target 972};
    \node[text=profilequery,font=\bfseries] at (12.35,3.50) {query \(x\)};
    \draw[gray!55] (10.55,4.18) -- (14.12,6.68);
    \draw[gray!55] (10.55,4.38) -- (13.98,6.78);
    \draw[dashed,draw=profileposition] (12.65,3.78) -- (12.65,7.08);
    \draw[dashed,draw=profilematch] (10.15,5.28) -- (14.55,5.28);
    \node[text=profileposition,anchor=south,fill=white,inner sep=1pt]
      at (12.65,3.82) {\(i=151\), \texttt{MHWVRQA}};
    \node[text=profilematch,anchor=east,fill=white,inner sep=1pt]
      at (14.47,5.45) {\(j=149\), \texttt{MHWVREA}};
    \node[fill=black,rounded corners=0.5mm,minimum width=3.2mm,
          minimum height=1.2mm,inner sep=0pt] at (11.85,5.10) {};
    \node[fill=black,rounded corners=0.5mm,minimum width=3.2mm,
          minimum height=1.2mm,inner sep=0pt] at (12.72,5.72) {};
    \node[fill=gray!30,draw=gray!60,rounded corners=0.5mm,
          minimum width=3.2mm,minimum height=1.2mm,inner sep=0pt]
      at (13.48,6.25) {};
    \node[align=center,fill=white,inner sep=1pt] at (12.32,6.72)
      {diagonal \(d=j-i=-2\)};
    \draw[->,draw=profilehit,line width=0.85pt]
      (8.38,4.70) -- (10.15,4.70);

    \node[draw=profilehit,rounded corners=2pt,anchor=south west,
          text width=13.65cm,align=left,inner sep=2.2mm]
      at (0.48,0.30) {
        \textcolor{profilequery}{for each query \(Q\)}\\
        \quad\textcolor{profileposition}{for each position \(i\) in \(Q\)}\\
        \qquad\textcolor{profilehit}{for each indexed hit \((t,j)\) of a
          similar \(k\)-mer}\\
        \qquad\quad \(d\gets j-i\)\\
        \qquad\quad if \(\operatorname{lastDiagonal}[t]=d\), emit candidate
          \((t,d)\)\\
        \qquad\quad \(\operatorname{lastDiagonal}[t]\gets d\)
      };
  \end{tikzpicture}
  \caption{Profile-search candidate reduction based on the staged organization
    of MMseqs2 \cite{mmseqs2_2017}. Panel (a) uses two successive \(k\)-mer
    hits on one diagonal to select candidates for ungapped scoring before final
    profile alignment. Panel (b) traverses queries, query positions, and index
    postings. The coordinate example illustrates the previous-diagonal check.
    Candidate reductions are qualitative.}
  \label{fig:profile-double-hit}
\end{figure}

For a hit at query position \(i\) and target position \(j\), the diagonal is
\(d=j-i\). The index and previous-diagonal state avoid an exhaustive comparison
of every profile symbol with every query symbol.

\section{Linear and graph reference models}
\label{sec:graph-approaches}

A linear reference provides one coordinate order for chaining. A variation
graph can represent shared sequence through common nodes and branch where
alternatives differ. A query alignment then chooses both a path through the
graph and an edit path against that path \cite{garrison_variation_graph_2018}.
A de Bruijn
graph represents local \(k\)-mer adjacency, commonly with \((k-1)\)-mers as
nodes and observed \(k\)-mers as edges \cite{pevzner_tang_waterman_2001}. It is
useful for assembly, but a graph walk is not itself an alignment.

Dynamic programming over an acyclic sequence graph can follow a topological
node order \cite{lee_partial_order_2002}. Cycles, repeated nodes, and many
equivalent paths require further constraints or search. Graph mappers still
need candidate reduction because one seed may occur on many paths. A graph
position identifies a node, offset, and orientation. Path identifiers and
graph versions are also needed when results use persistent path coordinates.

BioSeqDB presently chains anchors in a linear reference coordinate system. A
graph-native route would require graph positions in the seed index and
reachability-aware chaining. It would also require graph alignment or local
subgraph extraction, plus stable path identifiers in SQL results. The graph
model is therefore a change to the search object, not a substitute for scoring,
seeding, or alignment.

\section{Limitations of existing tools}
\label{sec:limitations-existing-tools}

BioSeqDB addresses database integration limits rather than mapper efficiency.
Minimap2 already provides an optimized nucleotide mapper and reusable
reference indexes \cite{minimap2_paper}. BLAST+ and MMseqs2 support large
sequence-database searches \cite{blast_plus_2009,mmseqs2_2017}. The relevant
limits are the trade-offs of sparse search and the boundary between standalone
tools and relational execution.

\subsection{Limits of sparse alignment search}

Minimizers reduce work by retaining only part of the available exact evidence.
Sensitivity depends on \(k\), the window length, ambiguous symbols, the error
pattern, and occurrence filtering. A short or damaged query may contain no
useful seed. A repetitive seed may produce many locations and weak positional
evidence. No later alignment kernel can recover a locus removed at this stage.

Chaining is exact only for its score and admitted predecessor set. Distance
limits, bandwidth, occurrence thresholds, and early stopping can remove a
useful chain. Final dynamic programming is likewise exact only over its
evaluated cells and endpoints. A narrow band, an adaptive block, X-drop, or
Z-drop may exclude the unconstrained optimum. Repeats can produce several
similar mappings, so MAPQ depends on alternatives and repetitiveness rather
than one alignment score. These are properties to report, not defects that
database integration removes.

Processor choice adds another trade-off. SIMD works best when lanes perform
similar operations over regular memory. GPU execution adds submission,
transfer, synchronization, and residency costs that need a sufficiently large
batch. A faster kernel does not imply a faster query when candidate
construction, data conversion, or result materialization dominates.

\subsection{Limits of a file-oriented execution boundary}

A conventional workflow exports selected sequences, invokes a mapper, parses
its output, and joins the records back to metadata. File formats remain useful
for interchange and archiving, but as a mandatory query boundary they hide
alignment work from the relational planner, repeat conversion, and leave index
compatibility to orchestration code. PostgreSQL functions can express typed
one-to-many computation, but a custom scan is needed to own batching, index
reuse, rescans, and result projection \cite{postgresql17_custom_scan}.

Representation and lifetime are part of this problem. Packed, out-of-line,
fragmented, and lazy sequences have different access costs. PostgreSQL may need
to detoast a variable-length value before access \cite{postgresql17_toast}.
Eager conversion can process a complete sequence when search needs one region.
A reusable index separately needs source and parameter identity, safe
publication, invalidation, and bounded loading. These are database lifecycle
requirements rather than changes to minimizer or alignment arithmetic.
